% !TEX root = main.tex

In a preliminary set of experiments, we study the behavior of \OURS during training.
We then qualitatively assess the filters learned with \OURS before comparing our approach to previous state-of-the-art models on standard benchmarks.

\begin{figure}[t]
	\centering
	\subfloat[Clustering quality]{
    \includegraphics[width=0.32\linewidth]{figures-figure_ratio_edge.pdf}
		\label{fig:edges}
	}
	\subfloat[Cluster reassignment]{
    \includegraphics[width=0.32\linewidth]{figures-figure_reassign.pdf}
		\label{fig:reassing}
	}
	\subfloat[Influence of k]{
    \includegraphics[width=0.32\linewidth]{figures-figure_perf_k.pdf}
		\label{fig:choice-k}
	}
	\caption{
		Preliminary studies.
		\protect\subref{fig:edges}: evolution of the clustering quality along training epochs;
		\protect\subref{fig:reassing}: evolution of cluster reassignments at each clustering step;
		\protect\subref{fig:choice-k}: validation mAP classification performance for various choices of $k$.
	}
	\label{fig:prelim}
\end{figure}

%%%%%%%%%%%%%%%%%%%%%%%%%%%%%%%
%%% Preliminary experiments %%%
%%%%%%%%%%%%%%%%%%%%%%%%%%%%%%%

\begin{figure}[t]
  \centering
  \begin{tabular}{cc}
    \includegraphics[width=0.48\linewidth]{images-filters-rgb-large.jpg}&
    \includegraphics[width=0.48\linewidth]{images-filters-sob-large.jpg}
  \end{tabular}
  \caption{Filters from the first layer of an AlexNet trained on unsupervised ImageNet on raw RGB input (left) or after a Sobel filtering (right).
  }
  \label{fig:filters}
\end{figure}

\subsection{Preliminary study}

We measure the information shared between two different assignments $A$ and $B$ of the same data by the
Normalized Mutual Information (NMI), defined as:
\begin{equation*}
\mathrm{NMI}(A;B)=\frac{\mathrm{I}(A;B)}{\sqrt{\mathrm{H}(A) \mathrm{H}(B)}}
\end{equation*}
where $\mathrm{I}$ denotes the mutual information and $\mathrm{H}$ the entropy. This measure can be applied to any assignment coming
from the clusters or the true labels.
If the two assignments $A$ and $B$ are independent, the NMI is equal to $0$.
If one of them is deterministically predictable from the other, the NMI is equal to $1$.
\\

\noindent\textbf{Relation between clusters and labels.}
Figure~\subref*{fig:edges} shows the evolution of the NMI between the cluster assignments and the ImageNet labels during training.
It measures the capability of the model to predict class level information.
Note that we only use this measure for this analysis and not in any model selection process.
The dependence between the clusters and the labels increases over time, showing that our features progressively capture information
related to object classes.
\\

\noindent\textbf{Number of reassignments between epochs.}
At each epoch, we reassign the images to a new set of clusters, with no guarantee of stability.
Measuring the NMI between the clusters at epoch $t-1$ and $t$ gives an insight on the actual
stability of our model.
Figure~\subref*{fig:reassing} shows the evolution of this measure during training.
The NMI is increasing, meaning that there are less and less reassignments and the clusters are stabilizing over time.
However, NMI saturates below $0.8$, meaning that a significant fraction of images are regularly reassigned between epochs.
In practice, this has no impact on the training and the models do not diverge.
\\

\noindent\textbf{Choosing the number of clusters.}
We measure the impact of the number $k$ of clusters used in $k$-means on the quality of the model.
We report the same down-stream task as in the hyperparameter selection process, \ie mAP on the \textsc{Pascal} VOC 2007 classification validation set.
We vary $k$ on a logarithmic scale, and report results after $300$ epochs in Figure~\subref*{fig:choice-k}.
The performance after the same number of epochs for every $k$ may not be directly comparable, but it reflects the hyper-parameter selection process
used in this work.
The best performance is obtained with $k=10,000$.
Given that we train our model on ImageNet, one would expect $k = 1000$ to yield the best results, but apparently some amount of over-segmentation is beneficial.

% The last experiment that we describe in this prelimnary study, is measuring the amount of reassignments that happen along training.
%%%%%%%%%%%%%%%%%%%%%%
%%% Visualizations %%%
%%%%%%%%%%%%%%%%%%%%%%

\subsection{Visualizations}
\label{sec:viz}

\noindent\textbf{First layer filters.}
Figure~\ref{fig:filters} shows the filters from the first layer of an AlexNet trained with \OURS on raw RGB images and images preprocessed with a Sobel filtering.
The difficulty of learning convnets on raw images has been noted before~\cite{bojanowski2017unsupervised,doersch2015unsupervised,noroozi2016unsupervised,paulin2015local}.
As shown in the left panel of Fig.~\ref{fig:filters}, most filters capture only color information that typically plays a little role for object classification~\cite{van2011evaluating}.
Filters obtained with Sobel preprocessing act like edge detectors.
\\

\begin{figure}[t]
\centering
\begin{tabular}{cccccccc}
  \multicolumn{2}{c}{\texttt{conv1}} &~~~& \multicolumn{2}{c}{\texttt{conv3}} &~~~& \multicolumn{2}{c}{\texttt{conv5}}\\
  \includegraphics[width=0.15\linewidth]{images-filters-smlayer0-channel43.jpeg}&
  \includegraphics[width=0.15\linewidth]{images-filters-layer0-channel43.jpeg}&&
  \includegraphics[width=0.15\linewidth]{images-filters-smlayer2-channel5.jpeg}&
  \includegraphics[width=0.15\linewidth]{images-filters-layer2-channel5.jpeg}&&
  \includegraphics[width=0.15\linewidth]{images-filters-smlayer4-channel45.jpeg}&
  \includegraphics[width=0.15\linewidth]{images-filters-layer4-channel45.jpeg}
\\
  \includegraphics[width=0.15\linewidth]{images-filters-smlayer0-channel46.jpeg}&
  \includegraphics[width=0.15\linewidth]{images-filters-layer0-channel46.jpeg}&&
  \includegraphics[width=0.15\linewidth]{images-filters-smlayer2-channel80.jpeg}&
  \includegraphics[width=0.15\linewidth]{images-filters-layer2-channel80.jpeg}&&
  \includegraphics[width=0.15\linewidth]{images-filters-smlayer4-channel2.jpeg}&
  \includegraphics[width=0.15\linewidth]{images-filters-layer4-channel2.jpeg}
\end{tabular}
\caption{
  Filter visualization and top $9$ activated images from a subset of $1$ million images from YFCC100M for target filters in the
  layers \texttt{conv1}, \texttt{conv3} and \texttt{conv5} of an AlexNet trained with \OURS on ImageNet.
  The filter visualization is obtained by learning an input image that maximizes the response to a target filter~\cite{yosinski2015understanding}.
}
\label{fig:activ}
\end{figure}

\begin{figure}[t]
\centering
\begin{tabular}{ccccccc}
Filter $0$ && Filter $33$ && Filter $145$ && Filter $194$
\\
%{\itshape Piano} && {\itshape Stone arcades} && {\itshape Car} && {\itshape Dog}
%\\
\includegraphics[width=0.24\linewidth]{images-filters-layer4-channel0.jpeg}&&
\includegraphics[width=0.24\linewidth]{images-filters-layer4-channel33.jpeg}&&
\includegraphics[width=0.24\linewidth]{images-filters-layer4-channel145.jpeg}&&
\includegraphics[width=0.24\linewidth]{images-filters-layer4-channel194.jpeg}
\\
Filter $97$ && Filter $116$ && Filter $119$ && Filter $182$
\\
%{\itshape Abstract shape} && {\itshape Drawing} && {\itshape Depth of field} && {\itshape Background blur}
%\\
\includegraphics[width=0.24\linewidth]{images-filters-layer4-channel97.jpeg}&&
\includegraphics[width=0.24\linewidth]{images-filters-layer4-channel116.jpeg}&&
\includegraphics[width=0.24\linewidth]{images-filters-layer4-channel119.jpeg}&&
\includegraphics[width=0.24\linewidth]{images-filters-layer4-channel182.jpeg}
\end{tabular}
\caption{
  Top $9$ activated images from a random subset of $10$ millions images from YFCC100M for target filters in the last convolutional layer.
  The top row corresponds to filters sensitive to activations by images containing objects.
  The bottom row exhibits filters more sensitive to stylistic effects.
  For instance, the filters $119$ and $182$ seem to be respectively excited by background blur and depth of field effects.
}
\label{fig:waouh}
\end{figure}

\noindent\textbf{Probing deeper layers.}
We assess the quality of a target filter by learning an input image that maximizes its activation~\cite{erhan2009visualizing,zeiler2014visualizing}.
We follow the process described by Yosinki~\etal~\cite{yosinski2015understanding} with a
cross entropy function between the target filter and the other filters of the same layer.
%Starting from Gaussian noise, we iteratively update the input to maximize the response of a target filter.
%applying on the image a Gaussian blur every $5$ updates and a $\ell_2$ penalization.
Figure~\ref{fig:activ} shows these synthetic images as well as the $9$ top activated images from a subset of $1$ million images
from YFCC100M.
As expected, deeper layers in the network seem to capture larger textural structures.
However, some filters in the last convolutional layers seem to be simply replicating
the texture already captured in previous layers, as shown on the second row of Fig.~\ref{fig:waouh}.
This result corroborates the observation by Zhang~\etal~\cite{zhang2016split} that features
from \texttt{conv3} or \texttt{conv4} are more discriminative than those from \texttt{conv5}.

Finally, Figure~\ref{fig:waouh} shows the top $9$ activated images of some \texttt{conv5} filters that seem to be semantically coherent.
The filters on the top row contain information about structures that highly corrolate with object classes.
The filters on the bottom row seem to trigger on style, like drawings or abstract shapes.


%%%%%%%%%%%%%%%%%%%%%%%%%%
%%% Linear Classifiers %%%
%%%%%%%%%%%%%%%%%%%%%%%%%%

\subsection{Linear classification on activations}
\label{sec:linear}

Following Zhang~\etal~\cite{zhang2016split}, we train a linear classifier on top of different frozen convolutional layers.
%This experiment exhibits at which layer the difference between supervised and unsupervised features.
This layer by layer comparison with supervised features exhibits where a convnet starts to be task specific, \ie specialized in object classification.
%To this end, we learn linear classifiers on top of the convolutional layers of an AlexNet exactly following the experimental setting of Zhang~\etal~\cite{zhang2016split}.
We report the results of this experiment on ImageNet and the Places dataset~\cite{zhou2014learning} in Table~\ref{tab:linear}.
We choose the hyperparameters by cross-validation on the training set.
On ImageNet, \OURS outperforms the state of the art from \texttt{conv3} to \texttt{conv5} layers by $3-5\%$.
%\footnote{All percentages indicated in comparisons are absolute, not relative, ie. the difference between 10\% and 12\% is 2\% not 20\%.}
The largest improvement is observed in the \texttt{conv4} layer, while the \texttt{conv1} layer performs poorly,
probably because the Sobel filtering discards color.
%This loss of accuracy at low-level layers can be related to the use of a Sobel filtering of the input, discarding color information.
Consistently with the filter visualizations of Sec.~\ref{sec:viz}, \texttt{conv3} works better than \texttt{conv5}.
Finally, the difference of performance between \OURS and a supervised AlexNet grows significantly on higher layers:
at layers \texttt{conv2-conv3} the difference is only around $6\%$, % bridging the gap between supervised and unsupervised approaches.
but this difference rises to $14.4\%$ at \texttt{conv5}, marking where the AlexNet probably stores most of the class level information.
In the supplementary material, we also report the accuracy if a MLP is trained on the last layer; \OURS outperforms the state of the art by $8\%$.

\begin{table}[t]
  \centering
  \resizebox{\columnwidth}{!}{%
    \begin{tabular}{@{}l c ccccc c ccccc@{}}
      \toprule
            &~~~& \multicolumn{5}{c}{ImageNet} &~~~& \multicolumn{5}{c}{Places} \\
            \cmidrule{3-7} \cmidrule{9-13}
      Method && \texttt{conv1} & \texttt{conv2} & \texttt{conv3} & \texttt{conv4} & \texttt{conv5} && \texttt{conv1} & \texttt{conv2} & \texttt{conv3} & \texttt{conv4} & \texttt{conv5} \\
      \midrule
      Places labels                                     && --   & --   & --   & --   & --             && $22.1$ & $35.1$ & $40.2$ & $43.3$ & $44.6$  \\
      ImageNet labels                                   && $19.3$ & $36.3$ & $44.2$ & $48.3$ & $50.5$           && $22.7$ & $34.8$ & $38.4$ & $39.4$ & $38.7$  \\
      %ImageNet labels Sobel                            && $11.3$ & $29.6$ & $39.2$ & $38.2$ & $48.5$           && $16.9$ & $31.7$ & $39.1$ & $41.2$ & $39.3$  \\
      Random                                            && $11.6$ & $17.1$ & $16.9$ & $16.3$ & $14.1$           && $15.7$ & $20.3$ & $19.8$ & $19.1$ & $17.5$  \\
      \midrule
      Pathak~\etal~\cite{pathak2016context}             && $14.1$ & $20.7$ & $21.0$ & $19.8$ & $15.5$           && $18.2$ & $23.2$ & $23.4$ & $21.9$ & $18.4$  \\
      Doersch~\etal~\cite{doersch2015unsupervised}      && $16.2$ & $23.3$ & $30.2$ & $31.7$ & $29.6$           && $19.7$ & $26.7$ & $31.9$ & $32.7$ & $30.9$  \\
      Zhang~\etal~\cite{zhang2016colorful}              && $12.5$ & $24.5$ & $30.4$ & $31.5$ & $30.3$           && $16.0$ & $25.7$ & $29.6$ & $30.3$ & $29.7$  \\
      Donahue~\etal~\cite{donahue2016adversarial}       && $17.7$ & $24.5$ & $31.0$ & $29.9$ & $28.0$           && $21.4$ & $26.2$ & $27.1$ & $26.1$ & $24.0$  \\
      Noroozi and Favaro~\cite{noroozi2016unsupervised} && $\textbf{18.2}$ & $28.8$ & $34.0$ & $33.9$ & $27.1$  && $23.0$ & $32.1$ & $35.5$ & $34.8$ & $31.3$  \\
      Noroozi~\etal~\cite{noroozi2017representation}    && $18.0$ & $\textbf{30.6}$ & $34.3$ & $32.5$ & $25.7$           && $\textbf{23.3}$ & $\textbf{33.9}$ & $36.3$ & $34.7$ & $29.6$  \\
      Zhang~\etal~\cite{zhang2016split}                 && $17.7$ & $29.3$ & $35.4$ & $35.2$ & $32.8$           && $21.3$ & $30.7$ & $34.0$ & $34.1$ & $32.5$  \\
      \midrule
      \OURS  && $12.9$ & $29.2$ & $\textbf{38.2}$ & $\textbf{39.8}$ & $\textbf{36.1}$ && $18.6$ & $30.8$ & $\textbf{37.0}$ & $\textbf{37.5}$ & $\textbf{33.1}$ \\
      \bottomrule
    \end{tabular}
  }
  \vspace{\soustable}
  \caption{
    Linear classification on ImageNet and Places using activations from the convolutional layers of an AlexNet as features.
    We report classification accuracy on the central crop.
    % Noroozi and Favaro~\cite{noroozi2016unsupervised} use stride $2$ instead of $4$ in the \texttt{conv1} layer.
    Numbers for other methods are from Zhang~\etal~\cite{zhang2016split}.
  }
  \label{tab:linear}
\end{table}

The same experiment on the Places dataset provides some interesting insights:
%Again, \OURS outperforms the state-of-the-art for most layers.
like \OURS, a supervised model trained on ImageNet suffers from a decrease of performance for higher layers (\texttt{conv4} versus \texttt{conv5}).
Moreover, \OURS yields \texttt{conv3-4} features that are comparable to those trained with ImageNet labels.
This suggests that when the target task is sufficently far from the domain covered by ImageNet, labels are less important.


%%%%%%%%%%%%%%%%%%%%%%
%%% Transfer tasks %%%
%%%%%%%%%%%%%%%%%%%%%%

\subsection{\textsc{Pascal} VOC 2007}
\label{sec:pascal}
Finally, we do a quantitative evaluation of \OURS on image classification, object detection and semantic segmentation on \textsc{Pascal} VOC.
The relatively small size of the training sets on \textsc{Pascal} VOC ($2,500$ images) makes this setup closer to a ``real-world'' application,
where a model trained with heavy computational resources, is adapted to a task or a dataset with a small number of instances.
Detection results are obtained using \texttt{fast-rcnn}\footnote{https://github.com/rbgirshick/py-faster-rcnn}; segmentation results are obtained using the code of Shelhamer~\etal\footnote{https://github.com/shelhamer/fcn.berkeleyvision.org}.
For classification and detection, we report the performance on the test set of \textsc{Pascal} VOC 2007 and choose our hyperparameters on the validation set.
For semantic segmentation, following the related work, we report the performance on the validation set of \textsc{Pascal} VOC 2012.

\begin{table}[t]
  \centering
  \begin{tabular}{lc cc c cc c cc}
    \toprule
    &~~~~~~~~& \multicolumn{2}{c}{Classification} &~~~& \multicolumn{2}{c}{Detection} &~~~& \multicolumn{2}{c}{Segmentation} \\
                      \cmidrule{3-4} \cmidrule{6-7} \cmidrule{9-10}
    Method && \textsc{fc6-8}& \textsc{all} && \textsc{fc6-8}& \textsc{all} && \textsc{fc6-8} & \textsc{all} \\
    \midrule
    ImageNet labels                                             && $~78.9~$ & $~79.9~$ && --       & $~56.8~$ && --       & $~48.0~$ \\
    Random-rgb                                                  && $~33.2~$ & $~57.0~$ && $~22.2~$ & $~44.5~$ && $~15.2~$ & $~30.1~$ \\
    Random-sobel                                                && $~29.0~$ & $~61.9~$ && $~18.9~$ & $~47.9~$ && $~13.0~$ & $~32.0~$ \\
    \midrule
    Pathak~\etal~\cite{pathak2016context}                       && $~34.6~$ & $~56.5~$ && --   & $~44.5~$ && --   & $~29.7~$ \\
    Donahue~\etal~\cite{donahue2016adversarial}$^*$             && $~52.3~$ & $~60.1~$ && --   & $~46.9~$ && --   & $~35.2~$ \\
    Pathak~\etal~\cite{pathak2016learning}                      && --       & $~61.0~$ && --   & $~52.2~$ && --   & --     \\
		Owens~\etal~\cite{owens2016ambient}$^*$                     && $~52.3~$ & $~61.3~$ && --   & --   && --   & --   \\
    Wang and Gupta~\cite{wang2015unsupervised}$^*$              && $~55.6~$ & $~63.1~$ && $~32.8^\dagger$   & $~47.2~$ && $~26.0^\dagger$   & $~35.4^\dagger$   \\
    Doersch~\etal~\cite{doersch2015unsupervised}$^*$            && $~55.1~$ & $~65.3~$ && --   & $~51.1~$ && --   & --   \\
    Bojanowski and Joulin~\cite{bojanowski2017unsupervised}$^*$ && $~56.7~$ & $~65.3~$ && $~33.7^\dagger$ & $~49.4~$ && $~26.7^\dagger$ & $~37.1^\dagger$ \\ % our detection-all 49.6
    Zhang~\etal~\cite{zhang2016colorful}$^*$                    && $~61.5~$ & $~65.9~$ && $~43.4^\dagger$ & $~46.9~$ && $~35.8^\dagger$ & $~35.6~$ \\ % our detection-all 47.6 // our sementation-all 38.8
    Zhang~\etal~\cite{zhang2016split}$^*$                       && $~63.0~$ & $~67.1~$ && --   & $~46.7~$ && --   & $~36.0~$ \\
    Noroozi and Favaro~\cite{noroozi2016unsupervised}           && --       & $~67.6~$ && --   & $~53.2~$ && --   & $~37.6~$  \\
    Noroozi~\etal~\cite{noroozi2017representation}              && --       & $~67.7~$ && --   & $~51.4~$ && --   & $~36.6~$ \\
		\midrule
    \OURS && $\textbf{70.4}$ & $\textbf{73.7}$ && $\textbf{51.4}$ & $\textbf{55.4}$ && $\textbf{43.2}$ & $\textbf{45.1}$ \\
    \bottomrule
  \end{tabular}
  \vspace{\soustable}
  \caption{
    Comparison of the proposed approach to state-of-the-art unsupervised feature learning on classification, detection and segmentation on \textsc{Pascal} VOC.
    $^*$ indicates the use of the data-dependent initialization of Kr\"ahenb\"uhl~\etal~\cite{krahenbuhl2015data}.
    Numbers for other methods produced by us are marked with a $\dagger$.
  }
  \label{tab:voc}
\end{table}

Table~\ref{tab:voc} summarized the comparisons of \OURS with other feature-learning approaches on the three tasks.
As for the previous experiments, we outperform previous unsupervised methods on all three tasks, in every setting.
The improvement with fine-tuning over the state of the art is the largest on semantic segmentation ($7.5\%$).
On detection, \OURS performs only slightly better than previously published methods.
Interestingly, a fine-tuned random network performs comparatively to many unsupervised methods, but performs poorly if only \textsc{fc6-8} are learned.
For this reason, we also report detection and segmentation with \textsc{fc6-8} for \OURS and a few baselines.
These tasks are closer to a real application where fine-tuning is not possible.
It is in this setting that the gap between our approach and the state of the art is the greater (up to $9\%$ on classification).

%Interestingly, on the classification task, the difference with the state-of-the-art is larger when finetuning the fully-connected layers ($7.9\%$) than when finetuning the whole network ($5.2\%$).
%This suggest that reporting the performance of fully fine-tuned networks (\texttt{all}) may not be conclusive.
%Moreover, it is worth noting that the performance of \OURS with \textsc{fc6-8} performs better ($70.9\%$) than unsupervised approaches with \texttt{all} ($67.7\%$).
